\section{Results}
\label{sec:results}

\subsection{Compression Ratio Evolution}
The repository history shows a clear progression from a simple order-0 baseline toward stronger BWT-based statistical compressors. The key milestones are \texttt{v5}, which established the core statistical pipeline; \texttt{v9}, which refines \texttt{v5} into the current production version; \texttt{v6}, which pushed ratio further with per-block search; and \texttt{v8}, which optimized for speed.

\subsection{Benchmark Comparison}
\begin{table}[htbp]
\caption{Aggregate benchmark results from the 8-file repository corpus. Lower ratio and bits/symbol are better.}
\centering
\scriptsize
\begin{tabular}{lccccc}
\toprule
Tool & Comp. bytes & Avg. ratio & Bits/symbol & Avg. comp. time & Avg. decomp. time \\
\midrule
v5 & 7,194,380 & 54.77\% & 4.3814 & 98 ms & 101 ms \\
v6 & 7,145,565 & 54.32\% & 4.3516 & 789 ms & 126 ms \\
v7 & 7,818,121 & 59.52\% & 4.7612 & \textbf{24 ms} & \textbf{15 ms} \\
v8 & 10,038,825 & 76.42\% & 6.1136 & 26 ms & 24 ms \\
\textbf{v9} & 7,194,380 & 54.77\% & 4.3814 & 67 ms & 64 ms \\
v10 & 5,145,570 & \textbf{39.17\%} & \textbf{3.1336} & 530 ms & 82 ms \\
gzip & 7,814,799 & 59.49\% & 4.7592 & 111 ms & 22 ms \\
bzip2 & 7,209,293 & 54.88\% & 4.3905 & 177 ms & 90 ms \\
xz & 7,114,088 & 54.16\% & 4.3325 & 456 ms & 65 ms \\
zstd & 7,695,793 & 58.58\% & 4.6867 & 20 ms & 13 ms \\
\bottomrule
\end{tabular}
\label{tab:compression_comparison}
\end{table}

Table~\ref{tab:compression_comparison} shows four distinct optima. \texttt{v10} achieves the best overall compression ratio at 39.17\%, driven by its ability to compress file D from 2\,MB to 14 bytes. On the remaining seven files \texttt{v10} matches \texttt{v6}'s 54.32\%, which already beats bzip2's 54.88\%, although it still remains behind xz at 54.16\% on those files. \texttt{v8} is the clear speed winner among the project versions. \texttt{v9} is the refined \texttt{v5} production pipeline, offering the best speed/ratio balance for general use. \texttt{v10} builds on \texttt{v6}'s pipeline and is strictly better: identical output on all non-PRNG files, plus the ability to compress PRNG data that every other tool stores raw.

At the file level, \texttt{v10} matches \texttt{v6} byte-for-byte on all files except D, because \texttt{v10} is \texttt{v6}'s pipeline with PRNG detection added on top. Among project versions, \texttt{v10}/\texttt{v6} win files A, B, C, E, and G; bzip2 wins F; and xz wins H. On file D, \texttt{v10} wins outright with its 14-byte PRNG output. The contrast with \texttt{v8} is sharp: on files like B and G, the statistical variants compress to roughly half the size of \texttt{v8}.

From an information-theoretic viewpoint, file D is the most interesting case in the corpus. Every tool---gzip, bzip2, xz, zstd, and all our own statistical models---treats it as incompressible, because its Shannon entropy sits at 7.999 bits/byte, essentially the theoretical maximum. But file D is not actually random. It turns out to be the output of glibc's \texttt{rand()} function seeded with 1. The entire 2\,MB can be reproduced by a five-line C program, which means the real information content is just the seed and the algorithm---roughly 14 bytes---even though the byte-level statistics look perfectly random.

\texttt{v10} takes advantage of this. When entropy is too high for statistical compression, it tries to match the input against glibc \texttt{rand()} output for seeds 0--65535. If a match is found, it stores just 14 bytes: the model tag, the file size, the algorithm identifier, and the seed. The decompressor regenerates the file by re-running the generator. The result on file D is a 0.0007\% compression ratio, or about 142,857 to 1. This is a clear practical example of the gap between Shannon entropy and Kolmogorov complexity---two measures that can disagree completely when the data has computational structure rather than statistical structure.

Files F and H instead expose model mismatch. F is only weakly compressible and fairly binary-heavy, so the current statistical line captures less exploitable structure than bzip2, which explains why bzip2 reaches 81.06\% while \texttt{v9} and \texttt{v6} remain at 82.42\% and 81.70\%. H is clearly compressible, but xz reduces it to 2.87 bits/symbol whereas the project variants stay near 3.5 bits/symbol, indicating that substantial longer-range structure remains uncaptured by the present models.

\subsection{Performance Analysis}
From a systems perspective, stronger modeling pays off on structured inputs but the search cost grows quickly. Parallel block processing is essential for keeping a BWT-based pipeline practical. \texttt{v9} remains the best overall balance: close to \texttt{v6}'s ratio at much lower cost, and far more compact than \texttt{v8}.

\subsection{Version Evolution}
The repository evolved through nine major iterations, exploring different points in the compression design space:

\begin{table}[htbp]
\centering
\scriptsize
\begin{tabular}{llllll}
\toprule
\textbf{Version} & \textbf{Key Innovation} & \textbf{Avg Ratio} & \textbf{Change} & \textbf{Comp. Time} & \textbf{Status} \\
\midrule
v1 & Order-0 + Arithmetic & 71.44\% & --- & 95 ms & Baseline \\
v2 & Order-1 + Range coding & 62.74\% & $-$8.7pp & 2.0 s & Superseded \\
v3 & BWT preprocessing & 57.48\% & $-$5.3pp & 2.0 s & Superseded \\
v4 & libsais + AVX2 + parallel & 56.39\% & $-$1.1pp & 68 ms & Superseded \\
v5 & PPM Method C + PPMD escape & 54.77\% & $-$1.6pp & 73 ms & Superseded \\
v6 & Per-block multi-pipeline & 54.32\% & $-$0.45pp & 543 ms & Best stat. ratio \\
v7 & 2-way interleaved rANS \cite{giesen2014} & 59.2\% & +4.4pp & \textbf{22 ms} & Speed-focused \\
v8 & LZ77 hash-table matching & 76.42\% & +17.2pp & \textbf{14 ms} & Ultra-fast \\
v9 & Simplified v5 + ZRLE fix & 54.77\% & 0pp & 67 ms & \textbf{Production} \\
v10 & PRNG detection (Kolmogorov) & \textbf{39.17\%} & $-$15.6pp & 530 ms & \textbf{Best ratio} \\
\bottomrule
\end{tabular}
\caption{Compression ratio evolution across project versions (lower ratio and time are better).}
\label{tab:version_evolution}
\end{table}

\subsection{Failed Experiment: Context Mixing}
Between \texttt{v5} and \texttt{v7}, the project explored a PAQ-inspired multi-order PPM experiment with context mixing. The approach used static weight training, train on the first 100KB, freeze weights, and store them in the header.

Results were decisive: 59.85\% average ratio (4.93pp worse than \texttt{v5}'s 54.73\%) while running 157$\times$ slower. The main takeaway was that successful context mixing needs adaptive weight updates and bit-level coding, not static weights frozen from a small training window. This influenced subsequent designs: \texttt{v6} returned to the proven \texttt{v5} core but added per-block transform search instead of model-order complexity.
